% !TEX program = xelatex
% =====================================================================
%  《复数》高三专题 · 课堂板书版 PPT —— 例题与公式源码
%  例题全部取自《复数专题.pdf》，原题号见每题上方注释。
%  编译： xelatex 复数专题_公式源码.tex
% =====================================================================
\documentclass[12pt,a4paper]{article}
\usepackage{amsmath,amssymb}
\usepackage{xeCJK}
\usepackage[margin=2.4cm]{geometry}
\setCJKmainfont{Microsoft YaHei}
\setmainfont{Times New Roman}
\linespread{1.3}
\setlength{\parindent}{0pt}
\newcommand{\dsmtitle}[1]{\par\medskip\noindent{\large\bfseries #1}\par\smallskip}
\newcommand{\dsmsource}[1]{\par\smallskip\noindent{\small\itshape 来源：#1}\par\medskip}
\begin{document}
\begin{center}
{\LARGE\bfseries 复数}\\[4pt]
{\small 高三数学专题 · 八种题型例题（LaTeX 公式源码）}
\end{center}
\vspace{0.4em}\hrule\vspace{0.8em}

% ---------- 原卷《复数专题》第 26 题 ----------
\dsmtitle{题型一｜复数的四则运算}
\noindent\textbf{例题}\quad
已知$i$是虚数单位，化简\par
\[\frac{2+3i}{1-2i}\]
的结果为$\underline{\hspace{4em}}$。\par
\dsmsource{天津市九校联考 2024--2025 学年高三下学期第一次模拟数学试卷（第 10 题）}

% ---------- 原卷《复数专题》第 39 题 ----------
\dsmtitle{题型二｜实部、虚部的求解}
\noindent\textbf{例题}\quad
已知$i$为虚数单位，$(3+2i)z=8+i$，\par
则$z$的虚部为$\underline{\hspace{4em}}$。\par
\dsmsource{天津市新华中学 2024--2025 学年高三下学期高考校模拟 2 数学试题（第 10 题）}

% ---------- 原卷《复数专题》第 40 题 ----------
\dsmtitle{题型三｜复数的模}
\noindent\textbf{例题}\quad
$i$是虚数单位，则\par
\[\left|\frac{3+4i}{1+i}\right|=\underline{\hspace{4em}}\]
\dsmsource{天津市武清区杨村第一中学 2025 届高三下学期第一次阶段性检测数学试卷（第 10 题）}

% ---------- 原卷《复数专题》第 36 题 ----------
\dsmtitle{题型四｜纯虚数、实数与参数}
\noindent\textbf{例题}\quad
已知复数\par
\[(a+3i)(1+2i)\]
是纯虚数，则实数$a$的值为$\underline{\hspace{4em}}$。\par
\dsmsource{《复数专题.pdf》原卷第 36 题}

% ---------- 原卷《复数专题》第 137 题 ----------
\dsmtitle{题型五｜复平面与象限}
\noindent\textbf{例题}\quad
设复数\par
\[z=\frac{4+2i}{1-i}\]
则$z$在复平面内对应的点位于第$\underline{\hspace{4em}}$象限。\par
\dsmsource{2026·南开中学五月考}

% ---------- 原卷《复数专题》第 58 题 ----------
\dsmtitle{题型六｜共轭复数}
\noindent\textbf{例题}\quad
已知$i$是虚数单位，则复数\par
\[\frac{2+i}{1-i}\]
的共轭复数为$\underline{\hspace{4em}}$。\par
\dsmsource{《复数专题.pdf》原卷第 58 题}

% ---------- 原卷《复数专题》第 45 题 ----------
\dsmtitle{题型七｜复数方程}
\noindent\textbf{例题}\quad
$i$是虚数单位，复数$z$满足\par
\[\frac{z-2i}{z}=i\]
则$z=\underline{\hspace{4em}}$。\par
\dsmsource{天津市河西区 2024--2025 学年高三下学期总复习质量调查二（二模）数学试题（第 10 题）}

% ---------- 原卷《复数专题》第 121 题 ----------
\dsmtitle{题型八｜综合判断与几何最值}
\noindent\textbf{例题}\quad
如果复数$z$满足\par
\[|z+1-i|=2\]
那么$|z-2+i|$的最大值是$\underline{\hspace{4em}}$。\par
\dsmsource{《复数专题.pdf》原卷第 121 题}

\end{document}
